% コマ6: call 前の16バイト整列 — 一時値が奇数個のときだけ sp が8ずれる
% return n * fact(n - 1); の生成中、sp が16バイト境界のどこに居るかを4場面で追う
\documentclass[border=8pt]{standalone}
% 講義資料の図(materials/sessions/figures/*.tex)共通プリアンブル
% 各図の .tex から % 講義資料の図(materials/sessions/figures/*.tex)共通プリアンブル
% 各図の .tex から % 講義資料の図(materials/sessions/figures/*.tex)共通プリアンブル
% 各図の .tex から \input{figure-preamble} で読み込む。
% 配色・フォントは latex/session-template.tex と揃える。

% 日本語(LuaLaTeX + luatexja)。等幅セル内の日本語も和文フォントで出す
\usepackage{luatexja}
\usepackage{luatexja-fontspec}
\setmainfont{Noto Sans CJK JP}
\setmainjfont{Noto Sans CJK JP}
\setmonofont{Inconsolata}[Scale=MatchLowercase]
\setmonojfont{Noto Sans CJK JP}

\usepackage{xcolor}
\definecolor{TcAccentDark}{HTML}{583B66}
\definecolor{TcAccentLight}{HTML}{EEE6F2}
\definecolor{TcAccentLighter}{HTML}{F7F3F9}
\definecolor{TcEmph}{HTML}{E64B17}
\definecolor{TcText}{HTML}{262626}
\definecolor{TcGrayText}{HTML}{737373}
\definecolor{TcGrayBg}{HTML}{F2F2F2}

\usepackage{tikz}
\usetikzlibrary{arrows.meta,positioning,calc,decorations.pathreplacing,fit,backgrounds}

\tikzset{
  % テキスト
  every node/.style={text=TcText},
  annot/.style={font=\small, text=TcText, align=left},
  gannot/.style={font=\small, text=TcGrayText, align=center},
  addr/.style={font=\ttfamily\small, text=TcGrayText},
  % メモリセル
  memcell/.style={draw=TcAccentDark, line width=0.6pt, fill=white,
                  minimum height=8.5mm, minimum width=40mm,
                  font=\ttfamily, inner sep=2pt, outer sep=0pt},
  memcell gray/.style={memcell, fill=TcGrayBg},
  memcell hi/.style={memcell, fill=TcAccentLighter},
  memcell dashed/.style={memcell, dashed, fill=none, text=TcGrayText},
  % 領域(セクション図など)
  region/.style={draw=TcAccentDark, line width=0.6pt, fill=white,
                 minimum width=48mm, align=center, inner sep=3pt, outer sep=0pt},
  % 制御フロー図のボックス
  cfgbox/.style={draw=TcAccentDark, line width=0.6pt, fill=white, rounded corners=1.5mm,
                 minimum width=28mm, minimum height=8mm, font=\small, align=center,
                 inner sep=2pt},
  % 矢印
  ptr/.style={-{Stealth[length=2.8mm]}, line width=1.1pt, color=TcEmph},
  flow/.style={-{Stealth[length=2.4mm]}, line width=0.8pt, color=TcAccentDark},
  flow back/.style={flow, dashed},
  regptr/.style={-{Stealth[length=2.8mm]}, line width=1.1pt, color=TcAccentDark},
  % 補助
  bracestyle/.style={decorate, decoration={brace, amplitude=4pt}, line width=0.8pt,
                     color=TcAccentDark},
}
 で読み込む。
% 配色・フォントは latex/session-template.tex と揃える。

% 日本語(LuaLaTeX + luatexja)。等幅セル内の日本語も和文フォントで出す
\usepackage{luatexja}
\usepackage{luatexja-fontspec}
\setmainfont{Noto Sans CJK JP}
\setmainjfont{Noto Sans CJK JP}
\setmonofont{Inconsolata}[Scale=MatchLowercase]
\setmonojfont{Noto Sans CJK JP}

\usepackage{xcolor}
\definecolor{TcAccentDark}{HTML}{583B66}
\definecolor{TcAccentLight}{HTML}{EEE6F2}
\definecolor{TcAccentLighter}{HTML}{F7F3F9}
\definecolor{TcEmph}{HTML}{E64B17}
\definecolor{TcText}{HTML}{262626}
\definecolor{TcGrayText}{HTML}{737373}
\definecolor{TcGrayBg}{HTML}{F2F2F2}

\usepackage{tikz}
\usetikzlibrary{arrows.meta,positioning,calc,decorations.pathreplacing,fit,backgrounds}

\tikzset{
  % テキスト
  every node/.style={text=TcText},
  annot/.style={font=\small, text=TcText, align=left},
  gannot/.style={font=\small, text=TcGrayText, align=center},
  addr/.style={font=\ttfamily\small, text=TcGrayText},
  % メモリセル
  memcell/.style={draw=TcAccentDark, line width=0.6pt, fill=white,
                  minimum height=8.5mm, minimum width=40mm,
                  font=\ttfamily, inner sep=2pt, outer sep=0pt},
  memcell gray/.style={memcell, fill=TcGrayBg},
  memcell hi/.style={memcell, fill=TcAccentLighter},
  memcell dashed/.style={memcell, dashed, fill=none, text=TcGrayText},
  % 領域(セクション図など)
  region/.style={draw=TcAccentDark, line width=0.6pt, fill=white,
                 minimum width=48mm, align=center, inner sep=3pt, outer sep=0pt},
  % 制御フロー図のボックス
  cfgbox/.style={draw=TcAccentDark, line width=0.6pt, fill=white, rounded corners=1.5mm,
                 minimum width=28mm, minimum height=8mm, font=\small, align=center,
                 inner sep=2pt},
  % 矢印
  ptr/.style={-{Stealth[length=2.8mm]}, line width=1.1pt, color=TcEmph},
  flow/.style={-{Stealth[length=2.4mm]}, line width=0.8pt, color=TcAccentDark},
  flow back/.style={flow, dashed},
  regptr/.style={-{Stealth[length=2.8mm]}, line width=1.1pt, color=TcAccentDark},
  % 補助
  bracestyle/.style={decorate, decoration={brace, amplitude=4pt}, line width=0.8pt,
                     color=TcAccentDark},
}
 で読み込む。
% 配色・フォントは latex/session-template.tex と揃える。

% 日本語(LuaLaTeX + luatexja)。等幅セル内の日本語も和文フォントで出す
\usepackage{luatexja}
\usepackage{luatexja-fontspec}
\setmainfont{Noto Sans CJK JP}
\setmainjfont{Noto Sans CJK JP}
\setmonofont{Inconsolata}[Scale=MatchLowercase]
\setmonojfont{Noto Sans CJK JP}

\usepackage{xcolor}
\definecolor{TcAccentDark}{HTML}{583B66}
\definecolor{TcAccentLight}{HTML}{EEE6F2}
\definecolor{TcAccentLighter}{HTML}{F7F3F9}
\definecolor{TcEmph}{HTML}{E64B17}
\definecolor{TcText}{HTML}{262626}
\definecolor{TcGrayText}{HTML}{737373}
\definecolor{TcGrayBg}{HTML}{F2F2F2}

\usepackage{tikz}
\usetikzlibrary{arrows.meta,positioning,calc,decorations.pathreplacing,fit,backgrounds}

\tikzset{
  % テキスト
  every node/.style={text=TcText},
  annot/.style={font=\small, text=TcText, align=left},
  gannot/.style={font=\small, text=TcGrayText, align=center},
  addr/.style={font=\ttfamily\small, text=TcGrayText},
  % メモリセル
  memcell/.style={draw=TcAccentDark, line width=0.6pt, fill=white,
                  minimum height=8.5mm, minimum width=40mm,
                  font=\ttfamily, inner sep=2pt, outer sep=0pt},
  memcell gray/.style={memcell, fill=TcGrayBg},
  memcell hi/.style={memcell, fill=TcAccentLighter},
  memcell dashed/.style={memcell, dashed, fill=none, text=TcGrayText},
  % 領域(セクション図など)
  region/.style={draw=TcAccentDark, line width=0.6pt, fill=white,
                 minimum width=48mm, align=center, inner sep=3pt, outer sep=0pt},
  % 制御フロー図のボックス
  cfgbox/.style={draw=TcAccentDark, line width=0.6pt, fill=white, rounded corners=1.5mm,
                 minimum width=28mm, minimum height=8mm, font=\small, align=center,
                 inner sep=2pt},
  % 矢印
  ptr/.style={-{Stealth[length=2.8mm]}, line width=1.1pt, color=TcEmph},
  flow/.style={-{Stealth[length=2.4mm]}, line width=0.8pt, color=TcAccentDark},
  flow back/.style={flow, dashed},
  regptr/.style={-{Stealth[length=2.8mm]}, line width=1.1pt, color=TcAccentDark},
  % 補助
  bracestyle/.style={decorate, decoration={brace, amplitude=4pt}, line width=0.8pt,
                     color=TcAccentDark},
}

\begin{document}
\begin{tikzpicture}[node distance=0mm,
    scell/.style={memcell, minimum width=30mm, minimum height=8mm},
    padcell/.style={memcell gray, minimum width=30mm, minimum height=8mm},
    freecell/.style={memcell dashed, minimum width=30mm, minimum height=8mm},
    stepttl/.style={annot, anchor=south, align=center},
    spmark/.style={annot, font=\ttfamily\bfseries, anchor=east, fill=white, inner sep=1pt},
    state/.style={annot, anchor=north, align=center},
    insn/.style={font=\ttfamily\footnotesize, text=TcText, align=center, anchor=north}]

  \node[gannot] at (69mm, 24mm)
      {{\ttfamily return n * fact(n - 1);} を生成するあいだの sp と16バイト境界};
  \node[annot, text=TcEmph, anchor=south, align=center] at (92mm, 15mm)
      {規約: {\ttfamily call} を出す時点の sp は16の倍数};

  % ---- 背景の16バイト境界 ----
  \begin{scope}[on background layer]
    \draw[dashed, line width=0.4pt, color=TcGrayText] (-30mm, 0mm) -- (158mm, 0mm);
    \draw[dashed, line width=0.4pt, color=TcGrayText] (-30mm, -16mm) -- (158mm, -16mm);
  \end{scope}
  \node[gannot, anchor=west] at (159mm, 0mm) {16バイト境界};
  \node[gannot, anchor=west] at (159mm, -16mm) {16バイト境界};

  % ---- ① プロローグ直後 ----
  \node[stepttl] at (0mm, 9mm) {① プロローグ直後};
  \node[spmark] (sp1) at (-17mm, 0mm) {sp};
  \draw[regptr] (sp1.east) -- (-15mm, 0mm);
  \node[state] at (0mm, -19mm) {{\ttfamily depth = 0}\\境界の上};
  \node[insn] at (0mm, -31mm) {addi sp, sp, -32};

  % ---- ② 左辺 n を積む ----
  \node[stepttl] at (46mm, 9mm) {② 左辺 {\ttfamily n} を積む};
  \node[scell, fill=TcAccentLight] (n2) at (46mm, -4mm) {n の値};
  \node[spmark] (sp2) at (29mm, -8mm) {sp};
  \draw[regptr] (sp2.east) -- (31mm, -8mm);
  \draw[bracestyle, color=TcEmph] (62mm, -16mm) -- (62mm, -8mm)
      node[midway, anchor=west, xshift=3mm, annot, text=TcEmph] {8 ずれ};
  \node[state] at (46mm, -19mm) {{\ttfamily depth = 1}\\境界から 8 ずれ};
  \node[insn] at (46mm, -31mm) {addi sp, sp, -8\\sd a0, 0(sp)};

  % ---- ③ call の直前 ----
  \node[stepttl] at (92mm, 9mm) {③ {\ttfamily call fact} の直前};
  \node[scell, fill=TcAccentLight] (n3) at (92mm, -4mm) {n の値};
  \node[padcell] (pad3) at (92mm, -12mm) {詰め物 8バイト};
  \node[spmark] (sp3) at (75mm, -16mm) {sp};
  \draw[regptr] (sp3.east) -- (77mm, -16mm);
  \node[state] at (92mm, -19mm) {{\ttfamily depth = 1}\\詰め物で境界の上へ};
  \node[insn] at (92mm, -31mm) {addi sp, sp, -8\\call fact};

  % ---- ④ 戻ってきたあと ----
  \node[stepttl] at (138mm, 9mm) {④ 戻ってきたあと};
  \node[scell, fill=TcAccentLight] (n4) at (138mm, -4mm) {n の値};
  \node[freecell] (pad4) at (138mm, -12mm) {詰め物は返した};
  \node[spmark] (sp4) at (121mm, -8mm) {sp};
  \draw[regptr] (sp4.east) -- (123mm, -8mm);
  \node[state] at (138mm, -19mm) {{\ttfamily depth = 1}\\{\ttfamily a1} へ復元して {\ttfamily mul}};
  \node[insn] at (138mm, -31mm) {addi sp, sp, 8\\ld a1, 0(sp)\\addi sp, sp, 8\\mul a0, a1, a0};

  % ---- パネル間の矢印 ----
  \foreach \x in {0, 46, 92} {
    \draw[flow, color=TcGrayText] (\x mm + 18mm, 4mm) -- (\x mm + 28mm, 4mm);
  }

  % ---- 補足 ----
  \node[gannot, anchor=north, align=center] at (69mm, -52mm)
      {引数の push は {\ttfamily call} の時点では戻し済みなので、ずれを決めるのは
       呼び出しを囲む式が積んだ一時値の個数だけである\\
       詰め物を入れるのは {\ttfamily call} の前後だけで、同じ幅を戻すため
       積んである一時値の位置は変わらない};
\end{tikzpicture}
\end{document}
